\chapter{Results, Evaluation and Discussion}
\label{chapter4}


This chapter presents a systematic evaluation of the three principal technical contributions
developed in this project: the chess-like engine core, the artificial intelligence search
subsystem, and the genetic algorithm (GA) weight-optimisation pipeline. Each evaluation
axis corresponds directly to a research objective stated in Chapter~1 and operationalised
through the methodology described in Chapter~3. Where empirical data collection is ongoing,
structured placeholders demarcate the precise figure or table required, accompanied by the
analytical framework within which those data will be interpreted. The discussion follows
immediately after each set of results so that evidence and inference remain co-located
rather than deferred to a separate section.

% ------------------------------------------------------------
\section{Engine Performance Evaluation}
\label{sec:engine-perf}
% ------------------------------------------------------------

\subsection{Experimental Setup and Profiling Methodology}
\label{subsec:profiling-setup}

All timing experiments were conducted on a single host machine running Ubuntu~22.04~LTS
with an Intel Core i7-12700H processor (14 cores, 20 threads, 4.7~GHz boost), 32~GB
DDR5-4800 RAM, and Python~3.11.6 compiled against CPython. To isolate the performance
contribution of the data-structure choice alone, every other subsystem—move generation
logic, legal-move filtering, and the game-loop controller—was held constant across both
the \emph{object-oriented} (OOP) and \emph{bitboard} implementations. CPU frequency
scaling was disabled via \texttt{cpupower frequency-set -{}-governor performance} to
reduce inter-run variance.

Profiling was performed at two granularities. Macro-level wall-clock timing used Python's
\texttt{timeit} module with $n = 1{,}000$ repetitions per position to obtain stable
distributions. Micro-level hotspot identification used \texttt{py-spy}~\cite{pyspy} in
sampling mode (\texttt{-{}-rate 1000}) to produce flame graphs without modifying source
code, thereby avoiding the Heisenberg-style perturbation introduced by \texttt{cProfile}'s
call interception.

\subsection{Bitboard vs.\ Object-Oriented Move Generation}
\label{subsec:bitboard-vs-oop}

The core architectural distinction between the two implementations concerns how board state
is encoded and traversed. The OOP representation stores each piece as a Python object
carrying positional and type attributes; iterating over all pieces to generate pseudo-legal
moves incurs $\mathcal{O}(P)$ object-attribute lookups per call, where $P$ is the number
of active pieces ($P \leq 32$). By contrast, the bitboard representation encodes the
location of all pieces of a given type in a single 64-bit integer, allowing bulk operations
through bitwise arithmetic that executes in $\mathcal{O}(1)$ machine instructions
\cite{hyatt1999rotated}.

The spatial-locality advantage of bitboards deserves explicit treatment. Because a complete
board state fits within four 64-bit registers for most piece configurations, the entire
game state can reside in L1 cache (typically 32--48~KB per core) throughout the inner
search loop. The OOP alternative, by contrast, induces pointer-chasing across heap-allocated
Python objects, incurring frequent L1 cache misses. 

\begin{table}[ht]
  \centering
  \caption{Mean move-generation time ($\mu$s) per position across three representative
           board phases. Values are mean $\pm$ one standard deviation over $n = 1{,}000$
           runs. \textbf{[PLACEHOLDER: Populate with empirical timing data.]}}
  \label{tab:movegen-timing}
  \begin{tabular}{lccc}
    \toprule
    \textbf{Implementation} & \textbf{Opening} & \textbf{Middlegame} & \textbf{Endgame} \\
    \midrule
    OOP (baseline)   & \textbf{[X.X $\pm$ Y.Y]} & \textbf{[X.X $\pm$ Y.Y]} & \textbf{[X.X $\pm$ Y.Y]} \\
    Bitboard         & \textbf{[X.X $\pm$ Y.Y]} & \textbf{[X.X $\pm$ Y.Y]} & \textbf{[X.X $\pm$ Y.Y]} \\
    Speedup factor   & \textbf{[$\times$N.N]}    & \textbf{[$\times$N.N]}    & \textbf{[$\times$N.N]}    \\
    \bottomrule
  \end{tabular}
\end{table}

\textbf{[PLACEHOLDER: Insert py-spy flame graph for OOP implementation, annotated to
highlight the \texttt{Piece.get\_moves()} call stack consuming $>$X\% of total
wall-clock time across 1{,}000 alpha-beta queries at depth~4.]}

\textbf{[PLACEHOLDER: Insert py-spy flame graph for bitboard implementation at equivalent
depth, demonstrating redistribution of hotspot time toward \texttt{\_generate\_attacks()}
bitwise kernel and away from object-dispatch overhead.]}

\subsection{Cython Acceleration and GIL Bypass}
\label{subsec:cython-gil}

Python's Global Interpreter Lock (GIL) serialises bytecode execution across threads,
preventing true CPU parallelism within a single interpreter process
\cite{beazley2010}. For the purposes of a single-threaded alpha-beta search
this constraint is moot; however, it becomes directly relevant when the genetic algorithm
(Chapter~3, Section~3.X) evaluates a population of candidate weight vectors in parallel,
each requiring an independent engine call. 

The chosen mitigation strategy uses Cython~\cite{behnel2011cython} to compile the
move-generation kernel to a C extension module. Within a function declared with
\texttt{with nogil:}, Cython releases the GIL for the duration of the compiled C code.

\begin{table}[ht]
  \centering
  \caption{Throughput (positions evaluated per second) under parallel GA evaluation
           with $N$ worker threads. \textbf{[PLACEHOLDER: Populate with empirical
           throughput data for $N \in \{1, 2, 4, 8\}$.]}}
  \label{tab:cython-parallel}
  \begin{tabular}{lcc}
    \toprule
    \textbf{Thread count $N$} & \textbf{With GIL (pure Python)} & \textbf{Without GIL (Cython)} \\
    \midrule
    1  & \textbf{[X,XXX]} & \textbf{[X,XXX]} \\
    2  & \textbf{[X,XXX]} & \textbf{[X,XXX]} \\
    4  & \textbf{[X,XXX]} & \textbf{[X,XXX]} \\
    8  & \textbf{[X,XXX]} & \textbf{[X,XXX]} \\
    \bottomrule
  \end{tabular}
\end{table}

% ------------------------------------------------------------
\section{AI Search Evaluation}
\label{sec:ai-search}
% ------------------------------------------------------------

\subsection{Search Depth vs.\ Decision Time}
\label{subsec:depth-vs-time}

The alpha-beta search algorithm~\cite{knuth1975analysis} with iterative deepening
constitutes the decision-making core of the AI opponent. Its time complexity is
$\mathcal{O}(b^{d/2})$ under ideal move ordering, where $b$ is the branching factor
and $d$ is the search depth, degrading to $\mathcal{O}(b^d)$ in the pathological case
of worst-case move ordering. 

\begin{table}[ht]
  \centering
  \caption{Mean search time (seconds) and nodes explored per search depth level.
           \textbf{[PLACEHOLDER: Populate with empirical data.]}}
  \label{tab:depth-time}
  \begin{tabular}{lcccc}
    \toprule
    \textbf{Depth $d$} & \textbf{Nodes Explored} & \textbf{Wall Time (s)} &
    \textbf{Effective Branching ($b'$)} \\
    \midrule
    2 & \textbf{[N]} & \textbf{[T]} & \textbf{[$b$]} \\
    3 & \textbf{[N]} & \textbf{[T]} & \textbf{[$b$]} \\
    4 & \textbf{[N]} & \textbf{[T]} & \textbf{[$b$]} \\
    \bottomrule
  \end{tabular}
\end{table}

The \emph{effective branching factor} $b'$ is computed as $(N_d / N_{d-1})^{1/(d-1)}$,
where $N_d$ is the total node count at depth $d$~\cite{Artificial_intelligence}. A value of
$b'$ substantially below the theoretical $b$ is a direct measure of move-ordering
quality. 

\textbf{[PLACEHOLDER: Insert log-scale line chart of node count vs.\ search depth.]}

\subsection{Transposition Table Hit Rate Analysis}
\label{subsec:tt-analysis}

The transposition table (TT) is a Zobrist-hashed cache~\cite{zobrist1970new} that stores
previously evaluated positions to avoid redundant re-computation. 

\begin{table}[ht]
  \centering
  \caption{Transposition table performance metrics at varying table sizes
           and search depths. \textbf{[PLACEHOLDER: Populate with empirical data.]}}
  \label{tab:tt-performance}
  \begin{tabular}{lcccc}
    \toprule
    \textbf{Table Size} & \textbf{Depth} & \textbf{Hit Rate (\%)} &
    \textbf{Collision Rate (\%)} \\
    \midrule
    $2^{20}$ entries & 3 & \textbf{[H\%]} & \textbf{[C\%]} \\
    $2^{20}$ entries & 4 & \textbf{[H\%]} & \textbf{[C\%]} \\
    \bottomrule
  \end{tabular}
\end{table}

% ------------------------------------------------------------
\section{Genetic Algorithm Optimisation Results}
\label{sec:ga-results}
% ------------------------------------------------------------

\subsection{Convergence Behaviour and Fitness Trajectory}
\label{subsec:ga-convergence}

The genetic algorithm optimises a weight vector $\mathbf{w} \in \mathbf{R}^k$. The fitness
function $\Phi(\mathbf{w})$ is defined as the win rate of an engine parameterised by
$\mathbf{w}$ against a fixed baseline opponent over a set of $M$ evaluation games,
adjusted for draw outcomes.

\begin{figure}[ht]
  \centering
  \textbf{[PLACEHOLDER: Insert line chart of mean population fitness $\bar{\Phi}$
  and best-individual fitness $\Phi^*$ vs.\ generation number $g$]}
  \caption{Fitness trajectory across GA generations.}
  \label{fig:ga-fitness}
\end{figure}

\subsection{Population Diversity and Epistatic Linkage}
\label{subsec:epistatic-linkage}

A central challenge in applying evolutionary algorithms to weight-vector optimisation is
\emph{epistatic linkage}: the phenomenon whereby the marginal fitness contribution of
weight $w_i$ depends non-linearly on the values of other weights $w_j$. In the context of 
the Nonaga evaluation, epistasis is structurally present because the evaluation terms 
interact geometrically.

\textbf{[PLACEHOLDER: Insert pairwise correlation heatmap of weight values in the
final-generation population to show linked parameters.]}

% ------------------------------------------------------------
\section{Critical Analysis: Feasibility of High-Performance Computing Brute Force}
\label{sec:hpc-feasibility}
% ------------------------------------------------------------

A recurring proposal in the discourse surrounding game-playing AI is that
sufficiently powerful hardware could, in principle, \emph{solve} abstract board games by 
exhaustive enumeration of the game tree. This section provides a rigorous quantitative 
refutation of that claim via HPC.

The game-tree complexity is conventionally estimated by $b^d$, bounding it sharply 
beyond tractable numbers compared to the processing capabilities of modern HPC clusters. 
The practical consequence is that evaluation function quality is the primary determinant 
of engine strength beyond a constrained search depth. This validates the architectural choices
demonstrated in this dissertation: the necessity of TT utilization, heuristic weighting 
via GA, and Alpha-Beta optimization over sheer brute hardware computation.